\begin{abstract}
Spatial transcriptomics platforms differ in panel, resolution and chemistry, so serial sections of one tissue yield
cell embeddings that disagree across technologies. Most integration methods, including spatially aware ones,
re-embed all data jointly, which moves the reference and invalidates everything built on it. We present \frost, a
post-hoc method that maps cells from a source platform into a frozen reference embedding with fused
Gromov--Wasserstein transport. Building on USHER, \frost computes the structure matrices once rather than in every EM
iteration, fits a linear identity-regularised map and ends with closed-form moment matching. On SPATCH data (Xenium
5K reference; CosMx 6K, Visium HD FFPE, Visium HD FF and Stereo-seq sources; three tumour types; 12 pairs), \frost
ranks among the top three of 11 methods on seven summary metrics. It transfers labels more accurately than the
spatially aware GraphST (balanced accuracy 0.72 vs 0.37; exact Wilcoxon $p=0.0005$) with no significant difference
in mixing, whereas GraphST is worse than no alignment and displaces its own reference (neighbourhood drift 0.98;
zero for \frost by construction). Against the USHER reference implementation, \frost reaches a similar total score
from the other side of the mixing--conservation trade-off, with 25-fold fewer structure computations and about
7.5-fold less EM time. A global map fitted to distributions and labels does blur gene-level neighbourhoods: once
leakage is removed, imputation and neighbour-based deconvolution favour unaligned or locally corrected embeddings,
which defines the scope of the method.
\end{abstract}
